\documentclass[11pt,a4paper]{article}

% ── Codificación, fuentes y lenguaje ───────────────────────────────────────────
\usepackage[utf8]{inputenc}
\usepackage[T1]{fontenc}
\usepackage[default,scale=0.95]{sourcesanspro}
\renewcommand{\familydefault}{\sfdefault}
\usepackage[spanish,es-noshorthands,es-tabla]{babel}

% ── Geometría y espaciado ─────────────────────────────────────────────────────
\usepackage[top=2.2cm, bottom=2.5cm, left=2.2cm, right=2.2cm, headheight=15pt]{geometry}
\usepackage{setspace}
\setstretch{1.18}
\usepackage{parskip}

% ── Matemáticas y unidades ────────────────────────────────────────────────────
\usepackage{amsmath}
\usepackage{amssymb}
\usepackage{siunitx}

% ── Circuitos ─────────────────────────────────────────────────────────────────
\usepackage[european, straightvoltages]{circuitikz}

% ── Tablas ────────────────────────────────────────────────────────────────────
\usepackage{booktabs}
\usepackage{tabularx}
\usepackage{array}
\usepackage{longtable}
\usepackage{colortbl}
\usepackage{enumitem}

% ── Colores, cajas e iconos ───────────────────────────────────────────────────
\usepackage[dvipsnames]{xcolor}
\usepackage{tcolorbox}
\tcbuselibrary{skins, breakable}
\usepackage{fontawesome5}
\usepackage{graphicx}

% ── Listados de código (dark-mode infográfico) ────────────────────────────────
%   Estilo base: fondo oscuro con números de línea. Los bloques lstlisting se
%   envuelven automáticamente en una caja tcolorbox redondeada.
\usepackage{listings}

% ── Secciones con barra de color (infografía) ────────────────────────────────
\usepackage{titlesec}
\titleformat{\section}
  {\normalfont\LARGE\bfseries\color{azulNoche}}
  {\colorbox{cyanNeon}{\color{fondoOscuro}\thesection}}{0.7em}{}
  [\vspace{2pt}\color{cyanNeon}\rule{\linewidth}{1.8pt}]
\titlespacing*{\section}{0pt}{24pt}{10pt}
\titleformat{\subsection}
  {\normalfont\large\bfseries\color{azulNoche}}
  {\textcolor{cyanNeon}{\thesubsection}}{0.7em}{}
  [\vspace{1pt}\color{grisLinea}\rule{\linewidth}{0.6pt}]
\titlespacing*{\subsection}{0pt}{14pt}{6pt}
\titleformat{\subsubsection}
  {\normalfont\normalsize\bfseries\color{azulNoche}}
  {\textcolor{cyanNeon}{\thesubsubsection}}{0.7em}{}
\titlespacing*{\subsubsection}{0pt}{10pt}{4pt}

% ── Cabeceras y pies de página ────────────────────────────────────────────────
\usepackage{fancyhdr}
\usepackage{lastpage}
\pagestyle{fancy}
\fancyhf{}
\renewcommand{\headrulewidth}{0pt}
\renewcommand{\footrulewidth}{0pt}
\fancyfoot[C]{%
  \begin{minipage}{\linewidth}
    \vspace{2pt}
    \color{cyanNeon}\rule{\linewidth}{1.2pt}\\[2pt]
    \centering\color{grisTexto}\footnotesize
    Página \thepage\ de \pageref{LastPage}
  \end{minipage}%
}

% ── Hiperenlaces ──────────────────────────────────────────────────────────────
\usepackage[colorlinks=true, linkcolor=azulNoche, urlcolor=cyanNeon]{hyperref}
\sloppy
\emergencystretch 3em

% ── Paleta de colores — tecnológico dark-mode ──────────────────────────────────
\definecolor{fondoOscuro}{HTML}{0D1B2A}     % Dark navy — fondo banda título
\definecolor{verdeTurquesa}{HTML}{004D40}    % Verde turquesa — bandas de marca
\definecolor{azulNoche}{HTML}{1B3A6B}       % Midnight blue — secciones, marcos
\definecolor{azulMedio}{HTML}{2A5298}       % Azul medio — variante de acento
\definecolor{cyanNeon}{HTML}{00C8E0}        % Cyan eléctrico — acentos primarios
\definecolor{verdeSignal}{HTML}{00B887}     % Verde señal — fortalezas, éxito
\definecolor{naranjaVivo}{HTML}{FF7043}     % Naranja vivo — mejoras, advertencias
\definecolor{rojoAlerta}{HTML}{E53935}      % Rojo alerta — errores
\definecolor{grisPapel}{HTML}{F4F6FA}       % Fondo tarjetas claras
\definecolor{grisLinea}{HTML}{C8D0E0}       % Bordes sutiles
\definecolor{grisTexto}{HTML}{6B7A99}       % Texto secundario
\definecolor{amarilloNota}{HTML}{FFB300}    % Notas / advertencias doradas

% Colores específicos para bloques de código (dark-mode)
\definecolor{codigoFondo}{HTML}{1E2D3D}      % Fondo del bloque de código
\definecolor{codigoComentario}{HTML}{7A8FA6} % Comentarios
\definecolor{codigoCadena}{HTML}{FF9D5C}     % Cadenas / strings
\definecolor{codigoKeyword}{HTML}{00C8E0}    % Palabras clave
\definecolor{codigoNumero}{HTML}{00E5A0}     % Números

% ── Banda de título: portada infográfica ──────────────────────────────────────
%   Diseño en 2 capas: banda de color (por defecto fondo oscuro) + franja de
%   acento cyan abajo. Uso: \bandaTitulo[color]{título}{subtítulo}.
%   El icono se puede cambiar con \renewcommand{\iconoBanda}{...} (FontAwesome).
\newcommand{\iconoBanda}{microchip}
\newcommand{\bandaTitulo}[3][fondoOscuro]{%
  \begin{tcolorbox}[
    enhanced,
    colback=#1, colframe=#1,
    boxrule=0pt, arc=8pt, outer arc=8pt,
    width=\linewidth,
    top=18pt, bottom=6pt, left=14pt, right=14pt,
    borderline south={3.5pt}{0pt}{cyanNeon},
  ]
    \begin{minipage}[c]{0.07\linewidth}
      \centering
      \textcolor{cyanNeon}{\fontsize{30}{30}\selectfont\faIcon{\iconoBanda}}
    \end{minipage}%
    \hspace{8pt}%
    \begin{minipage}[c]{0.88\linewidth}
      {\fontsize{20}{24}\selectfont\bfseries\color{white}#2}\par\vspace{5pt}%
      \textcolor{cyanNeon!80!white}{\small\faIcon{angle-right}~#3}%
    \end{minipage}
    \vspace{4pt}
  \end{tcolorbox}%
  \vspace{8pt}%
}

% ── Tarjetas de datos (infografía) ────────────────────────────────────────────
\tcbset{
  tarjetaDato/.style={
    enhanced, breakable,
    arc=6pt, outer arc=6pt,
    colback=grisPapel, colframe=azulNoche,
    boxrule=1pt,
    fonttitle=\bfseries\small, coltitle=white,
    attach boxed title to top left={yshift=-3mm, xshift=6mm},
    boxed title style={
      colback=azulNoche, colframe=azulNoche,
      arc=4pt, boxrule=0pt,
    },
    drop fuzzy shadow=azulNoche!30!white,
    top=8pt, bottom=8pt, left=10pt, right=10pt,
  },
  tarjetaIcono/.style={
    enhanced, breakable,
    arc=6pt, outer arc=6pt,
    colback=white, colframe=grisLinea, boxrule=0.8pt,
    fonttitle=\bfseries\small, coltitle=azulNoche,
    borderline west={3pt}{0pt}{cyanNeon},
    drop fuzzy shadow=azulNoche!25!white,
    top=6pt, bottom=6pt, left=10pt, right=10pt,
  },
  cajaTitulo/.style={
    enhanced, breakable,
    arc=6pt, outer arc=6pt,
    colback=azulNoche!10!grisPapel, colframe=azulNoche,
    boxrule=1pt,
    fonttitle=\bfseries, coltitle=white,
    attach boxed title to top left={yshift=-3mm, xshift=6mm},
    boxed title style={
      colback=azulNoche, colframe=azulNoche,
      arc=4pt, boxrule=0pt,
    },
    drop fuzzy shadow=azulNoche!25!white,
    top=8pt, bottom=8pt, left=10pt, right=10pt,
  },
  cajaContenido/.style={
    enhanced, breakable,
    arc=5pt, outer arc=5pt,
    colback=grisPapel, colframe=grisLinea,
    boxrule=0.8pt,
    fonttitle=\bfseries\small, coltitle=azulNoche,
    top=6pt, bottom=6pt, left=10pt, right=10pt,
  },
  cajaFortaleza/.style={
    enhanced, breakable,
    arc=6pt, outer arc=6pt,
    colback=verdeSignal!8!white, colframe=verdeSignal,
    boxrule=1pt,
    borderline west={4pt}{0pt}{verdeSignal},
    fonttitle=\bfseries\small, coltitle=verdeSignal!70!black,
    drop fuzzy shadow=verdeSignal!20!white,
    top=6pt, bottom=6pt, left=10pt, right=10pt,
  },
  cajaMejora/.style={
    enhanced, breakable,
    arc=6pt, outer arc=6pt,
    colback=naranjaVivo!8!white, colframe=naranjaVivo,
    boxrule=1pt,
    borderline west={4pt}{0pt}{naranjaVivo},
    fonttitle=\bfseries\small, coltitle=naranjaVivo!80!black,
    drop fuzzy shadow=naranjaVivo!20!white,
    top=6pt, bottom=6pt, left=10pt, right=10pt,
  },
  cajaRecomendacion/.style={
    enhanced, breakable,
    arc=6pt, outer arc=6pt,
    colback=cyanNeon!6!white, colframe=cyanNeon!60!azulNoche,
    boxrule=1pt,
    borderline west={4pt}{0pt}{cyanNeon},
    fonttitle=\bfseries\small, coltitle=azulNoche,
    drop fuzzy shadow=azulNoche!20!white,
    top=6pt, bottom=6pt, left=10pt, right=10pt,
  },
}

% ── Ícono + texto (encabezados de sección y elementos) ────────────────────────
\newcommand{\iconotexto}[2]{\textcolor{cyanNeon}{\faIcon{#1}}~#2}

% ── Listados de código: estilo dark + auto-envoltura ─────────────────────────
\lstdefinestyle{estiloCodigo}{
    backgroundcolor=\color{codigoFondo},
    basicstyle=\ttfamily\small\color{white},
    commentstyle=\color{codigoComentario}\itshape,
    stringstyle=\color{codigoCadena},
    keywordstyle=\color{codigoKeyword}\bfseries,
    numberstyle=\tiny\color{codigoComentario},
    numbers=left,
    numbersep=10pt,
    showstringspaces=false,
    breaklines=true,
    breakatwhitespace=false,
    tabsize=4,
    frame=none,
    captionpos=b,
    aboveskip=6pt,
    belowskip=4pt,
    xleftmargin=14pt,
    framexleftmargin=14pt,
    literate=
      {á}{{\'a}}1 {é}{{\'e}}1 {í}{{\'i}}1 {ó}{{\'o}}1 {ú}{{\'u}}1
      {Á}{{\'A}}1 {É}{{\'E}}1 {Í}{{\'I}}1 {Ó}{{\'O}}1 {Ú}{{\'U}}1
      {ñ}{{\~n}}1 {Ñ}{{\~N}}1 {ü}{{\"u}}1 {Ü}{{\"U}}1
      {¿}{{?`}}1 {¡}{{!`}}1
}
\lstdefinestyle{estiloPython}{
    style=estiloCodigo,
    language=Python,
}
\lstset{style=estiloCodigo}
\BeforeBeginEnvironment{lstlisting}{%
  \begin{tcolorbox}[
    enhanced, breakable,
    arc=6pt, outer arc=6pt,
    colback=codigoFondo, colframe=azulNoche,
    boxrule=1pt,
    drop fuzzy shadow=azulNoche!25!white,
    top=2pt, bottom=2pt, left=0pt, right=0pt,
  ]%
}
\AfterEndEnvironment{lstlisting}{\end{tcolorbox}}

% ── Cajas de contenido temáticas ──────────────────────────────────────────────
\newtcolorbox{cajaRecuerda}{
    enhanced, breakable,
    arc=6pt, outer arc=6pt,
    colback=azulNoche!10!grisPapel,
    colframe=azulNoche, boxrule=1pt,
    borderline west={4pt}{0pt}{cyanNeon},
    fonttitle=\bfseries\small\color{white},
    title={\faIcon{info-circle}~Recuerda},
    attach boxed title to top left={yshift=-3mm, xshift=6mm},
    boxed title style={colback=azulNoche, arc=4pt, boxrule=0pt},
    drop fuzzy shadow=azulNoche!25!white,
    left=10pt, right=10pt, top=8pt, bottom=8pt,
}

\newtcolorbox{cajaConcepto}{
    enhanced, breakable,
    arc=6pt, outer arc=6pt,
    colback=verdeSignal!8!white,
    colframe=verdeSignal, boxrule=1pt,
    borderline west={4pt}{0pt}{verdeSignal},
    fonttitle=\bfseries\small\color{white},
    title={\faIcon{lightbulb}~Concepto clave},
    attach boxed title to top left={yshift=-3mm, xshift=6mm},
    boxed title style={colback=verdeSignal!80!black, arc=4pt, boxrule=0pt},
    drop fuzzy shadow=verdeSignal!20!white,
    left=10pt, right=10pt, top=8pt, bottom=8pt,
}

\newtcolorbox{cajaEjemplo}{
    enhanced, breakable,
    arc=6pt, outer arc=6pt,
    colback=cyanNeon!5!white,
    colframe=cyanNeon!70!azulNoche, boxrule=1pt,
    borderline west={4pt}{0pt}{cyanNeon},
    fonttitle=\bfseries\small\color{fondoOscuro},
    title={\faIcon{code}~Ejemplo},
    attach boxed title to top left={yshift=-3mm, xshift=6mm},
    boxed title style={colback=cyanNeon, arc=4pt, boxrule=0pt},
    drop fuzzy shadow=azulNoche!20!white,
    left=10pt, right=10pt, top=8pt, bottom=8pt,
}

\newtcolorbox{cajaEjercicio}{
    enhanced, breakable,
    arc=6pt, outer arc=6pt,
    colback=azulMedio!7!white,
    colframe=azulMedio, boxrule=1pt,
    borderline west={4pt}{0pt}{azulMedio},
    fonttitle=\bfseries\small\color{white},
    title={\faIcon{pencil-alt}~Ejercicio},
    attach boxed title to top left={yshift=-3mm, xshift=6mm},
    boxed title style={colback=azulMedio, arc=4pt, boxrule=0pt},
    drop fuzzy shadow=azulNoche!20!white,
    left=10pt, right=10pt, top=8pt, bottom=8pt,
}

\newtcolorbox{cajaReto}{
    enhanced, breakable,
    arc=6pt, outer arc=6pt,
    colback=naranjaVivo!6!white,
    colframe=naranjaVivo, boxrule=1pt,
    borderline west={4pt}{0pt}{naranjaVivo},
    fonttitle=\bfseries\small\color{white},
    title={\faIcon{fire}~Reto},
    attach boxed title to top left={yshift=-3mm, xshift=6mm},
    boxed title style={colback=naranjaVivo, arc=4pt, boxrule=0pt},
    drop fuzzy shadow=naranjaVivo!20!white,
    left=10pt, right=10pt, top=8pt, bottom=8pt,
}

% ── Indicadores de dificultad ─────────────────────────────────────────────────
\newcommand{\facil}{\textcolor{verdeSignal}{\faIcon{circle}\,\textbf{Fácil}}}
\newcommand{\intermedio}{\textcolor{naranjaVivo}{\faIcon{circle}\,\textbf{Intermedio}}}
\newcommand{\dificil}{\textcolor{rojoAlerta}{\faIcon{circle}\,\textbf{Difícil}}}

% ─────────────────────────────────────────────────────────────────────────────

% ── Cabecera específica del reporte del skill ──
\renewcommand{\iconoBanda}{book}
\fancyhead[L]{\color{azulNoche}\small\bfseries \faIcon{book}~\textcolor{cyanNeon}{circuitos\_dispositivos\_electronicos}}
\fancyhead[R]{\color{grisTexto}\smallElectrónica: dispositivos semiconductores y circuitos analógicos}

\begin{document}
\bandaTitulo[fondoOscuro]{Skill: circuitos\_dispositivos\_electronicos}{Skill de referencia rápida para análisis de circuitos con dispositivos semiconductores (diodos, BJT, MOS, A.O.) y circuitos analógicos. Úsese cuando se requieran leyes de Kirchhoff, Thévenin/Norton, respuesta transitoria}

\section*{Circuitos y Dispositivos Electrónicos}


\subsection*{Cuándo usar este skill}


\begin{itemize}[leftmargin=*,topsep=2pt,itemsep=1pt]
  \item Resolver circuitos resistivos, RC, RL y con fuentes dependientes.
  \item Analizar diodos rectificadores, zener, LED, fotodiodos y circuitos recortadores/fijadores.
  \item Analizar y diseñar etapas con transistores bipolares (BJT) y MOS en continua y pequeña señal.
  \item Trabajar con amplificadores operacionales en región lineal y no lineal (comparadores, Schmitt trigger).
  \item Calcular equivalentes Thévenin/Norton, superposición, transferencia de señal y potencia.
  \item Interpretar modelos SPICE y hojas de especificaciones de dispositivos.
  \item Comprender fundamentos físicos: unión PN, efecto transistor, efecto de campo, capacidades parásitas.
\end{itemize}


\subsection*{Conceptos clave}


\begin{itemize}[leftmargin=*,topsep=2pt,itemsep=1pt]
  \item \textbf{Tensión, corriente y potencia}: definiciones, convenios de signos, analogía hidráulica.
  \item \textbf{Leyes de Kirchhoff}: LCK (nudos) y LTK (mallas).
  \item \textbf{Señales}: escalón, pulso, rampa, exponencial, sinusoidal; valor medio y eficaz; representación compleja.
  \item \textbf{Elementos ideales}: conductor, interruptor, fuentes independientes y dependientes.
  \item \textbf{Resistencia}: Ley de Ohm, conductancia, efecto Joule, asociaciones serie/paralelo, divisores de tensión/corriente.
  \item \textbf{Métodos de análisis}: nudos, mallas, superposición, Thévenin/Norton, fuente de prueba.
  \item \textbf{Condensador y bobina}: relaciones i-v, continuidad, asociaciones, constante de tiempo, respuesta transitoria y sinusoidal, impedancia, energía almacenada.
  \item \textbf{Transformador ideal y real}: relación de transformación, adaptación de impedancias.
  \item \textbf{Diodo}: modelo ideal, exponencial, por tramos lineales; recta de carga; aplicaciones (rectificador, detector de envolvente, recortador, fijador, protección).
  \item \textbf{Transistor bipolar (BJT)}: regiones de operación, βF, αF, efecto Early, modelos de pequeña señal (híbrido en π, parámetros h), etapas EC/CC/BC, par diferencial, CMRR, puertas lógicas.
  \item \textbf{Transistor MOS}: tipos (acumulación/vaciamiento, canal N/P), regiones (corte, óhmica, saturación), efecto sustrato, modulación de longitud de canal, inversor CMOS, puertas lógicas, amplificación, modelo de pequeña señal.
  \item \textbf{Amplificador operacional}: modelo ideal, cortocircuito virtual, configuraciones básicas, comparadores, realimentación.
  \item \textbf{Dispositivos de potencia}: SCR, triac, diac, GTO, IGBT, convertidor DC-DC reductor.
  \item \textbf{Optoelectrónica}: LED, fotodiodo, célula solar, fototransistor, optoacoplador.
  \item \textbf{Fundamentos físicos}: semiconductores intrínsecos/extrínsecos, unión PN, capacidades de transición y difusión, principio de operación de BJT y MOS.
  \item \textbf{SPICE}: directivas .OP, .TF, .TRAN, .AC, .PRINT, .PROBE; modelos de diodo, BJT y MOS.
\end{itemize}


\subsection*{Fórmulas esenciales}


\begin{itemize}[leftmargin=*,topsep=2pt,itemsep=1pt]
  \item Ley de Ohm: $v = iR$, $i = Gv$, $P_R = i^2R = v^2/R$.
  \item Divisor de tensión: $v_o = v_g \frac{R_2}{R_1+R_2}$.
  \item Divisor de corriente: $i_2 = i_g \frac{R_1}{R_1+R_2}$.
  \item Serie: $R_s = \sum R_i$; Paralelo: $\frac{1}{R_p} = \sum \frac{1}{R_i}$.
  \item Condensador: $i_c = C \frac{dv_c}{dt}$, $v_c(t) = v_c(0) + \frac{1}{C}\int_0^t i_c d\tau$, $Z_c = \frac{1}{j\omega C}$.
  \item Bobina: $v_L = L \frac{di_L}{dt}$, $i_L(t) = i_L(0) + \frac{1}{L}\int_0^t v_L d\tau$, $Z_L = j\omega L$.
  \item Constante de tiempo: $\tau = RC$ o $\tau = L/R$.
  \item Respuesta RC escalón: $v_c(t) = V_a(1 - e^{-t/RC})$, $i_c(t) = \frac{V_a}{R}e^{-t/RC}$.
  \item Respuesta RL escalón: $i(t) = \frac{V_a}{R}(1 - e^{-tR/L})$, $v_L(t) = V_a e^{-tR/L}$.
  \item Energía: $W_C = \frac{1}{2}CV^2$, $W_L = \frac{1}{2}LI^2$.
  \item Transformador ideal: $\frac{v_2}{v_1} = n$, $\frac{i_2}{i_1} = -\frac{1}{n}$, $R_e = \frac{R_L}{n^2}$.
  \item Diodo exponencial: $I_d = I_s(e^{v_D/\eta V_T} - 1)$, $V_T \approx 25$ mV a 300 K.
  \item Diodo por tramos: $v_D = V_\gamma + i_d R_s$ (conduce), $i_d = 0$ (corte).
  \item Rizado rectificador: $\Delta V_o \approx \frac{V_M}{f C R_L}$.
  \item BJT activa: $i_C = \beta_F i_B$, $i_E = (\beta_F+1)i_B$, $\alpha_F = \frac{\beta_F}{\beta_F+1}$.
  \item BJT pequeña señal: $g_m = \frac{I_{CQ}}{V_T}$, $r_\pi = \frac{\beta_F}{g_m}$, $r_o = \frac{V_A}{I_{CQ}}$.
  \item MOS saturación: $I_D = \frac{KW}{2L}(v_{GS}-V_T)^2(1+\lambda v_{DS})$.
  \item MOS óhmica: $I_D = K\frac{W}{L}\left[(v_{GS}-V_T)v_{DS} - \frac{v_{DS}^2}{2}\right]$.
  \item MOS pequeña señal: $g_m = 2\sqrt{K\frac{W}{L}I_{DQ}}$, $r_{ds} = \frac{1}{\lambda I_{DQ}}$.
  \item A.O. inversor: $\frac{v_o}{v_g} = -\frac{R_F}{R_s}$; no inversor: $1+\frac{R_1}{R_2}$.
  \item CMRR: $CMRR = 20\log(\rho)$ dB.
  \item Unión PN: $V_{bi} = V_T \ln\left(\frac{N_A N_D}{n_i^2}\right)$, $C_j = \frac{C_{jo}}{\sqrt{1-V/V_{bi}}}$.
  \item Capacidad de difusión: $C_s = \tau_t \frac{I_s}{V_T}e^{V/V_T}$.
\end{itemize}


\subsection*{Procedimiento básico de análisis}


\begin{enumerate}[leftmargin=*,topsep=2pt,itemsep=1pt]
  \item \textbf{Identificar} el tipo de circuito (resistivo, RC/RL, con diodos, con BJT/MOS, con A.O.).
  \item \textbf{Simplificar} si es posible: asociaciones serie/paralelo, equivalentes Thévenin/Norton.
  \item \textbf{Aplicar} leyes de Kirchhoff y ecuaciones de los dispositivos (hipótesis de región para no lineales).
  \item \textbf{Verificar} hipótesis (p. ej., diodo conduce/corta, BJT en activa/saturación, MOS en saturación/óhmica).
  \item \textbf{Calcular} variables de interés: tensiones, corrientes, potencias, ganancias, márgenes dinámicos.
  \item \textbf{Validar} con SPICE o comparación con resultados analíticos si es necesario.
\end{enumerate}


\subsection*{Tablas de referencia rápida}


\begin{itemize}[leftmargin=*,topsep=2pt,itemsep=1pt]
  \item Ver \texttt{tablas.md} para tablas de símbolos, prefijos, regiones de operación, configuraciones de A.O., etapas BJT, tipos de MOS, etc.
\end{itemize}


\subsection*{Glosario}


\begin{itemize}[leftmargin=*,topsep=2pt,itemsep=1pt]
  \item Ver \texttt{glosario.md} para definiciones de términos técnicos.
\end{itemize}


\subsection*{Fórmulas detalladas}


\begin{itemize}[leftmargin=*,topsep=2pt,itemsep=1pt]
  \item Ver \texttt{formulas.md} para el compendio completo de fórmulas.
\end{itemize}


\section*{Fórmulas de Electrónica}


\subsection*{Ley de Ohm y Potencia}


\begin{itemize}[leftmargin=*,topsep=2pt,itemsep=1pt]
  \item $v = i \cdot R$
  \item $p = v \cdot i = i^2 R = \frac{v^2}{R}$
\end{itemize}


\subsection*{Leyes de Kirchhoff}


\begin{itemize}[leftmargin=*,topsep=2pt,itemsep=1pt]
  \item KCL: $\sum i_{nudo} = 0$
  \item KVL: $\sum v_{malla} = 0$
\end{itemize}


\subsection*{Divisores}


\begin{itemize}[leftmargin=*,topsep=2pt,itemsep=1pt]
  \item Tensión: $v_o = v_{in} \frac{R_2}{R_1 + R_2}$
  \item Corriente: $i_2 = i_{in} \frac{R_1}{R_1 + R_2}$
\end{itemize}


\subsection*{Resistencias Equivalentes}


\begin{itemize}[leftmargin=*,topsep=2pt,itemsep=1pt]
  \item Serie: $R_{eq} = R_1 + R_2 + \dots$
  \item Paralelo: $\frac{1}{R_{eq}} = \frac{1}{R_1} + \frac{1}{R_2} + \dots$
\end{itemize}


\subsection*{Condensadores y Bobinas}


\begin{itemize}[leftmargin=*,topsep=2pt,itemsep=1pt]
  \item $i_c = C \frac{dv_c}{dt}$
  \item $v_L = L \frac{di_L}{dt}$
  \item Serie: $\frac{1}{C_{eq}} = \sum \frac{1}{C_i}$, $L_{eq} = \sum L_i$
  \item Paralelo: $C_{eq} = \sum C_i$, $\frac{1}{L_{eq}} = \sum \frac{1}{L_i}$
\end{itemize}


\subsection*{Transitorios RC/RL}


\begin{itemize}[leftmargin=*,topsep=2pt,itemsep=1pt]
  \item $\tau = RC$ o $\tau = L/R$
  \item $x(t) = X_f + (X_i - X_f)e^{-t/\tau}$
\end{itemize}


\subsection*{Impedancias (RPS)}


\begin{itemize}[leftmargin=*,topsep=2pt,itemsep=1pt]
  \item $Z_R = R$
  \item $Z_C = \frac{1}{j\omega C}$
  \item $Z_L = j\omega L$
\end{itemize}


\subsection*{Amplificador Operacional Ideal}


\begin{itemize}[leftmargin=*,topsep=2pt,itemsep=1pt]
  \item Inversor: $A_v = -\frac{R_F}{R_s}$
  \item No Inversor: $A_v = 1 + \frac{R_1}{R_2}$
  \item Seguidor: $A_v = 1$
  \item Sumador: $v_o = -R_F \sum \frac{v_i}{R_i}$
\end{itemize}


\subsection*{Diodo}


\begin{itemize}[leftmargin=*,topsep=2pt,itemsep=1pt]
  \item Shockley: $I_d = I_s (e^{v_D/(\eta V_T)} - 1)$
  \item $V_T \approx 25mV$ a 300K
\end{itemize}


\subsection*{BJT (Activa)}


\begin{itemize}[leftmargin=*,topsep=2pt,itemsep=1pt]
  \item $I_C = \beta_F I_B$
  \item $I_E = (\beta_F + 1)I_B$
  \item $g_m = I_{CQ} / V_T$
  \item $r_\pi = \beta_F / g_m$
\end{itemize}


\subsection*{MOSFET (Saturación)}


\begin{itemize}[leftmargin=*,topsep=2pt,itemsep=1pt]
  \item $I_D = \frac{1}{2} K \frac{W}{L} (V_{GS} - V_T)^2$
  \item $g_m = 2 \sqrt{K \frac{W}{L} I_{DQ}}$
\end{itemize}


\subsection*{Máxima Transferencia de Potencia}


\begin{itemize}[leftmargin=*,topsep=2pt,itemsep=1pt]
  \item $R_L = R_{th}$ para máxima potencia.
  \item $P_{max} = \frac{V_{th}^2}{4 R_{th}}$
\end{itemize}


\vspace{10pt}

\section*{Glosario de Términos}


\begin{itemize}[leftmargin=*,topsep=2pt,itemsep=1pt]
  \item \textbf{Ánodo (A):} Terminal positivo de un diodo.
  \item \textbf{Cátodo (K):} Terminal negativo de un diodo.
  \item \textbf{Circuito Abierto:} Interrupción en un circuito; corriente cero.
  \item \textbf{Cortocircuito:} Conexión directa entre dos puntos; tensión cero.
  \item \textbf{Cortocircuito Virtual:} Condición en A.O. ideal con realimentación negativa donde $v^+ \approx v^-$.
  \item \textbf{Diodo:} Dispositivo semiconductor que permite el paso de corriente en una sola dirección.
  \item \textbf{Diodo Zener:} Diodo diseñado para operar en la región de ruptura inversa, usado como referencia de tensión.
  \item \textbf{Efecto Early:} Modulación de la anchura de la base en un BJT que causa dependencia de $I_C$ con $V_{CE}$.
  \item \textbf{Efecto Sustrato:} Variación de la tensión umbral $V_T$ en un MOSFET debido a la tensión entre surtidor y sustrato.
  \item \textbf{Fuente Dependiente:} Fuente de tensión o corriente cuyo valor está controlado por otra tensión o corriente del circuito.
  \item \textbf{Fuente Independiente:} Fuente de tensión o corriente con valor fijo o definido por una función del tiempo.
  \item \textbf{LED (Diodo Emisor de Luz):} Diodo que emite luz cuando se polariza en directa.
  \item \textbf{Masa (Tierra):} Nodo de referencia en un circuito con potencial de 0V.
  \item \textbf{Nudo:} Punto de conexión entre dos o más componentes.
  \item \textbf{Malla:} Trayectoria cerrada en un circuito.
  \item \textbf{Realimentación Negativa:} Conexión de la salida a la entrada inversora de un amplificador, estabilizando la ganancia.
  \item \textbf{Rectificador:} Circuito que convierte corriente alterna (AC) en corriente continua (DC).
  \item \textbf{Señal:} Magnitud física (tensión o corriente) que varía en el tiempo para transportar información.
  \item \textbf{Tensión de Ruptura ($V_z$, BV):} Tensión inversa a la que un diodo comienza a conducir fuertemente.
  \item \textbf{Tensión Térmica ($V_T$):} $kT/q \approx 25mV$ a temperatura ambiente.
  \item \textbf{Tensión Umbral ($V_\gamma$, $V_T$):} Tensión mínima de polarización directa para que un diodo conduzca (\textasciitilde{}0.7V para Si), o tensión puerta-surtidor para formar el canal en un MOSFET.
  \item \textbf{Transconductancia ($g_m$):} Relación entre la variación de corriente de salida y la variación de tensión de entrada en un dispositivo activo.
\end{itemize}


\vspace{10pt}

\section*{Tablas de Referencia}


\subsection*{Unidades y Prefijos}


\begin{center}
\begin{tabular}{@{}lll@{}}
\toprule
Magnitud & Unidad & Símbolo \\
\midrule
Tensión & Voltio & V \\
Corriente & Amperio & A \\
Potencia & Vatio & W \\
Resistencia & Ohmio & $\Omega$ \\
Capacidad & Faradio & F \\
Inductancia & Henrio & H \\
\bottomrule
\end{tabular}
\end{center}



\begin{center}
\begin{tabular}{@{}lll@{}}
\toprule
Prefijo & Símbolo & Factor \\
\midrule
giga & G & $10^9$ \\
mega & M & $10^6$ \\
kilo & k & $10^3$ \\
mili & m & $10^{-3}$ \\
micro & $\mu$ & $10^{-6}$ \\
nano & n & $10^{-9}$ \\
pico & p & $10^{-12}$ \\
\bottomrule
\end{tabular}
\end{center}



\subsection*{Regiones de Operación del BJT (NPN)}


\begin{center}
\begin{tabular}{@{}llll@{}}
\toprule
Región & Unión BE & Unión BC & Característica \\
\midrule
Activa & Directa & Inversa & $I_C = \beta_F I_B$ \\
Corte & Inversa & Inversa & $I_C \approx 0$ \\
Saturación & Directa & Directa & $V_{CE} \approx 0.2V$ \\
\bottomrule
\end{tabular}
\end{center}



\subsection*{Regiones de Operación del MOSFET (Canal N)}


\begin{center}
\begin{tabular}{@{}lll@{}}
\toprule
Región & Condición & $I_D$ \\
\midrule
Corte & $V_{GS} < V_T$ & 0 \\
Óhmica & $V_{DS} < V_{GS} - V_T$ & $K \frac{W}{L}[(V_{GS}-V_T)V_{DS} - \frac{V_{DS}^2}{2}]$ \\
Saturación & $V_{DS} \ge V_{GS} - V_T$ & $\frac{1}{2} K \frac{W}{L} (V_{GS}-V_T)^2$ \\
\bottomrule
\end{tabular}
\end{center}



\subsection*{Configuraciones de Amplificadores (BJT)}


\begin{center}
\begin{tabular}{@{}llll@{}}
\toprule
Configuración & Ganancia de Tensión & Impedancia de Entrada & Impedancia de Salida \\
\midrule
Emisor Común & Alta (-) & Media & Media \\
Colector Común & $\approx 1$ & Alta & Baja \\
Base Común & Alta (+) & Baja & Alta \\
\bottomrule
\end{tabular}
\end{center}



\subsection*{Modelos Equivalentes}


\begin{center}
\begin{tabular}{@{}ll@{}}
\toprule
Teorema & Circuito Equivalente \\
\midrule
Thévenin & Fuente de tensión $V_{th}$ en serie con $R_{th}$ \\
Norton & Fuente de corriente $I_n$ en paralelo con $R_n$ \\
Relación & $V_{th} = I_n \cdot R_{th}$, $R_{th} = R_n$ \\
\bottomrule
\end{tabular}
\end{center}



\vspace{10pt}

\end{document}
